% SPDX-License-Identifier: GPL-3.0-or-later
% Catalan (català) -- what the reading editions' preamble leaves to the language.
% Input by lib/frank-preamble.tex as \input{\FrankLib/lang/\BookLang.tex};
% docs/languages.md section 6 lists the hooks every file here must define and
% section 12 the ones it may; the guide's page "Adding a language"
% (html-guide/markdown/lookup-and-languages/adding-a-language.md) walks
% through writing one.
%
% Scaffolded from lib/lang/_template.tex by lib/newlang.py.  Everything below
% the switches came out of lib/languages.json, so it agrees with the registry
% by construction; \FrankHowTo did not, and is the one thing you must write.
% The model to read while writing it is lib/lang/fa.tex.

% ---- the babel language and the switches ----------------------------------
% \FrankLangName is what \foreignlanguage and \begin{otherlanguage} are given
% everywhere in the core preamble: babel must know this name, or every passage
% comes out in the roman with a "Unknown language" error.
\newcommand{\FrankLangName}{catalan}
\FrankRTLfalse           % chunk left, gloss right; \pw needs no bidi isolate
\FrankHasBarefalse       % nothing comes off, so pass 3 would repeat pass 1
\FrankHasAltfalse        % whether \FrankAltFont exists (the title, the chapter number)
\FrankHasReadingfalse    % whether \chr's third argument is set as a reading
\FrankHasVertfalse       % whether pass 4 is \FrankVertPass rather than the font pass
\FrankHasAloudfalse      % whether the chunks carry words, and pass 2 is each one's reading alone

% ---- fonts ------------------------------------------------------------------
% No face is NAMED here: \FrankPickFont walks the registry's tex.main and then
% tex.main_fallbacks and takes the first one installed, so the .tex and the
% .json cannot drift apart (docs/languages.md section 12).  A bundled face --
% one that travels in lib/fonts/ -- is found because build.sh puts that
% directory on OSFONTDIR.
\babelprovide[import=ca]{catalan}
%% No font of its own: tex.main is null in the registry, so the text is
%% set in the roman the core preamble already loaded -- the same face as
%% the glosses, which is right for a Latin script and wrong for any other.

% ---- the vocabulary line -----------------------------------------------------
% \vb{infinitive}{sound}{1st sg. present}{sound}{past participle}{sound}{meaning}:
% these two are the labels printed before the second and third forms.  The
% three forms are the registry's "vb_forms", which the chunk editor names
% when somebody types a \vb:
%   "infinitive", "1st sg. present", "past participle"
% Whatever you choose here, docs/lang/ca.md must name the same three
% forms in the same order, or the annotator fills slots the page labels
% differently.
% The core prints the second and the third pair only when its FORM is given
% (lib/frank-preamble.tex), so a verb that has no such form leaves the pair
% blank, {}{}, rather than repeating another one.  What the three cannot say
% -- an auxiliary, a class, a governed case -- goes in ONE parenthesis after
% the meaning, items parted by "; ", each only where it is not the ordinary
% case: {to drive (er fährt; aux. sein)} is German's.
% Catalan gives the infinitive, the first person singular present and the
% masculine singular past participle: the present shows the stem's changes
% (serveixo, vinc, faig), and the participle is what every compound tense is
% built on (he parlat); the everyday past is va + the infinitive (va parlar),
% which the first form already gives.  No extras.
% The pair below comes from the registry's "vb_labels" ("pres.", "p.p."), which is
% where the reader takes it from as well: change one and change the other, or
% the PDF and the reader of one book label the same form differently.
% lib/newlang.py --check compares them.
\newcommand{\FrankVbPres}{pres.}
\newcommand{\FrankVbPast}{p.p.}

% ---- the two Lua filters ---------------------------------------------------
% frank_strip(s): the bare form of a chunk -- what pass 3 shows.  Written from
% the registry's `strip`; the identity when there is nothing to take off.
% frank_latin(s): the label's digits as ASCII, because a PDF destination name
% cannot carry the language's own figures.  Written from `digits`.
% Both are called for every language, so both must exist even when they do
% nothing.  No #, % or literal backslash in a \directlua body: TeX reads it
% first (NOTES section 9 of the Persian book).
\directlua{
  function frank_strip(s) return s end
  function frank_latin(s) return s end
}

% ---- the "How to read this" page -----------------------------------------------
% Printed by lib/frank-frontmatter.tex, before the first chapter: the only
% place the reader is told what the passes are for.  Write it in the second
% person, one paragraph per pass, then a paragraph on the transliteration
% scheme and one on how verbs are given -- the shape lib/lang/fa.tex has.
% \pw{...} sets one word of the language inside the English; it is defined by
% the core preamble, after this file, and only expanded when the page is set.
%
\newcommand{\FrankHowTo}{%
Every passage comes twice.

\smallskip
\textbf{First} the whole sentence, and nothing else. Try to read it. Get as
far as you can and do not worry about what you cannot get — the point is the
attempt, made before any help arrives.

\smallskip
\textbf{Second} the same sentence in small chunks, Catalan on the left, and
on the right three lines: how it sounds, then how to look each word up, then
what it means. Now you find out how much you had right. Catalan needs nothing
added to its letters to be read, so there is no third pass: the sentence you
attempted is already the sentence as it is printed.

\smallskip
The pronunciation line follows standard Central Catalan, the Catalan of
Barcelona, and it is given only where the spelling and a few rules do not tell
you the sound. The rules first, since the line does not repeat them: an
unstressed \pw{a} or \pw{e} is the neutral vowel \textit{ə}, the
\textit{a} of \textit{about} (\pw{casa}, \textit{kázə}), and an unstressed
\pw{o} is \textit{u}; the \pw{r} of an infinitive is not said
(\pw{parlar}, \textit{pərlá}), nor a \pw{t} after \pw{l} or \pw{n} at the end
of a word (\pw{molt}), nor the \pw{b} of \pw{amb}, nor \pw{h}; \pw{b} and
\pw{v} are one sound. What the line shows is the rest. The stressed vowel
carries an accent, and on an \textit{e} or an \textit{o} the accent says
which of two sounds it is — open \textit{è}, \textit{ò} against closed
\textit{é}, \textit{ó} — which Catalan spelling shows only where it writes
an accent anyway (\pw{porta}, \textit{pòrtə}; \pw{boca}, \textit{bókə}).
\textit{ʎ} is \pw{ll} and \textit{ɲ} is \pw{ny}; \textit{ʃ} is \pw{x} and
\pw{ix}, \textit{tʃ} is \pw{tx} and a final \pw{ig} (\pw{vaig},
\textit{vatʃ}), \textit{ʒ} is \pw{j} and \pw{g} before \textit{e} or
\textit{i}; a double \textit{ll} in the line is the long \textit{l} of
\pw{col·lecció}, written with the middle dot; \textit{rr} is the trill,
which \pw{r} also is at the start of a word. A book whose line is in IPA
writes \textit{ˈ} before the stressed syllable instead, and \textit{ɛ},
\textit{ɔ} for the open vowels. Valencian and Balearic texts are printed as
they are written, and their line follows the recording.

\smallskip
Verbs are given as infinitive, first person present, past participle,
meaning — the two later forms are the ones you cannot guess (\pw{tenir},
\pw{tinc}, \pw{tingut}; \pw{servir}, \pw{serveixo}, \pw{servit}), and a form
carries its sound beside it only where its stressed \textit{e} or \textit{o}
has no accent written on it. An irregular future stands in brackets after the
meaning (\textit{fut.} \pw{tindré}): the same stem makes the conditional,
\pw{tindria}. The past a story is told in is \pw{va} with the infinitive
— \pw{vaig parlar} is \textit{I spoke}, not \textit{I go to speak} — and it is
named after the verb's entry, as every form the entry does not show is. Catalan
has two verbs for \textit{to be}, and the line says which: \pw{ser} for what a
thing is and where it is, \pw{estar} for how it is at this moment. A
subjunctive is named where it appears, since your own language may have no
form for that mood.
}
